\subsection{Exact splitting and reconstruction of the faithful finite extension}

\begin{result}{A.4.19}{Proved}{The marked faithful extension is recovered from its module}
For $E=\mathcal P_c$ of A.4.18 let $A$ be its displayed
three-primary image of $T$.  There is the literal exact
sequence
\[
 1\longrightarrow A\cong C_3^{12}\longrightarrow E
   \xrightarrow{\ b\ }B=G_2\times H\longrightarrow1,
 \qquad
 G_2=Q_8\mathbin{\times}_{C_2^2}(\mathbb Z/4)^2,
 \quad H=C_3^2.                                         \tag{EX.1}
\]
The quotient map is the joint map induced by the
quaternion representation, ordinary abelian reduction
modulo four, and ordinary abelian reduction modulo three.
The representation of $B$ on $A$ factors through the
specified $Q_8\times H$, and $A^{G_2}=0$.
The extension is split.  Its splittings are conjugate by
unique elements of $A$.

The same statements hold for the characteristic common
source $E=\mathcal E^-$, with $A=C_3^{36}$ and exactly the
same quotient $B$.  In both cases the kernel is
characteristic: it is the intersection of the unique
normal Sylow-three subgroup with the kernel of the
maximal elementary abelian three-quotient.

After choosing one actual splitting $s:B\to E$, compatible
pairs $(f,\beta)$ of automorphisms of $A$ and $B$, defined
by
\[
 f\rho(b)f^{-1}=\rho(\beta(b))\qquad(b\in B),             \tag{EX.2}
\]
give the actual automorphisms
\[
       a\,s(b)\longmapsto f(a)\,s(\beta(b)).             \tag{EX.3}
\]
Every automorphism of the marked extension with the same
pair differs from (EX.3) by a unique inner automorphism
from $A$.  Thus (EX.3) gives the exact reconstruction
modulo that explicitly retained inner ambiguity.
For $q\in\Gamma(2)\bmod12$, the literal signed base map
(EX.5) has period matrix $q$ exactly when its displayed
$H$-label map is $q\bmod3$.  In the quaternion-centralizing
sector, passage to outer automorphisms further removes
the faithful inner action of $Z(Q_8)\times H$, of order
$18$, with identity period matrix.
\end{result}

\begin{proof}
The map to $G_2$ is the joint quaternion and modulo-four
map in (QS.7); its image has order $32$.  The modulo-three
map has image $H$ of order nine.  The joint image is their
product, since a common quotient of groups of these
coprime orders is trivial.  The order calculations in
A.4.18 therefore identify its kernel as the displayed
$A$: its order is $3^{12}$ for $\mathcal P_c$ and $3^{36}$
for $\mathcal E^-$.  The actual maps from $T$ identify
these kernels with the faithful modules already
constructed.  In particular the unique central
quaternion involution acts as $-I$.  Its fixed space
is zero, proving $A^{G_2}=0$.

Here is a finite group proof of splitting, including
its exact ambiguity.  Let $E_0=b^{-1}(G_2)$ and choose
a Sylow-two subgroup $P$ of $E_0$.  Its order is $32$,
so $b|_P:P\to G_2$ is an isomorphism.  For any $e\in E$,
the group $ePe^{-1}$ is another Sylow-two subgroup of
the normal subgroup $E_0$.  Sylow conjugacy gives
$ePe^{-1}=aPa^{-1}$ for some $a\in A$: indeed any
conjugating element in $E_0=AP$ can have its $P$
factor discarded.  Thus $a^{-1}e\in N_E(P)$ and
$N_E(P)$ maps onto all of $B$.
An element of $A$ normalizing $P$ centralizes it:
its commutators with $P$ lie in both $A$ and $P$.
Consequently
\[
 A\cap N_E(P)=C_A(P)=A^{G_2}=0.                         \tag{EX.4}
\]
The restriction $b:N_E(P)\to B$ is therefore an
isomorphism.  Its inverse is an actual splitting.

Every complement $S$ to $A$ contains a unique normal
Sylow-two subgroup $P_S$, because $S\cong G_2\times H$
and its Sylow-two factor is normal.  Hence
$S\subseteq N_E(P_S)$; both groups have order $|B|$,
so equality holds.  The preceding Sylow conjugacy
shows all complements are $A$-conjugate.  The
conjugating element is unique, by (EX.4).  This proves
both splitting and its ambiguity without deleting an
extension-class condition.

The quotient $E/A$ has maximal two-quotient $G_2$ and
maximal elementary abelian three-quotient $H$.
The preimage of $H$ is the unique normal Sylow-three
subgroup, of order $3^{14}$ or $3^{38}$ respectively.
The joint abelianization of $E$ is
$(\mathbb Z/12)^2$, by the exponent-lattice calculations
in A.4.18; therefore its maximal elementary abelian
three-quotient is exactly the displayed $H$.
The intersection in the statement is thus precisely
$A$ and is characteristic.

In the split coordinates the multiplication is
$(a,b)(a',b')=(a+\rho(b)a',bb')$.  Substitution proves
that (EX.3) is an automorphism exactly when (EX.2)
holds.  An automorphism inducing $(f,\beta)$ sends
the chosen complement to another complement; the
unique $A$-conjugation just proved corrects it to
(EX.3).  This proves the last assertion, including
its uniqueness.
\par\medskip\noindent\textbf{The exact signed base maps.}\quad
Write an element of $G_2$ as $(x,v)$ with
$x\in Q_8$, $v\in(\mathbb Z/4)^2$, and equal images
in $C_2^2$.  For every actual matrix
$q\in\Gamma(2)\bmod12$ and the specified label map
$\nu\in\operatorname{Aut}(H)$ there is the literal
base automorphism
\[
 (x,v,h)\longmapsto(x,(q\bmod4)v,\nu h).                \tag{EX.5}
\]
It is well-defined because $q\bmod2=I$.  On the
ordinary abelian quotient modulo twelve this has
matrix $q$ precisely when $\nu=q\bmod3$, in the
same downstairs basis.  Thus the equality between
these two label maps is an explicit compatibility
condition in the reconstruction, not an identification
of unrelated representations.

For automorphisms of $A$ centralizing the quaternion
action and normalizing $H$, (EX.2) reduces to their
specified $H$-label map and (EX.5).  Passing further
from automorphisms to outer automorphisms removes
the actual inner action of $Z(Q_8)\times H$, of
order $2\cdot9=18$, in this quaternion-centralizing
sector.  Its action on $A$ is faithful, and its
period matrix is the identity.  The other quaternion
inner actions are retained when the full quaternion
automorphism residue is restored.

\end{proof}
